\section{Die Maxwell-Gleichungen}
\begin{flushleft}

\subsection{Der Maxwell'sche Verschiebungsstrom}
\textbf{Lade-/ Entladevorgang:}
\[
    I = \frac{dQ}{dt} = \underbrace{\epsilon_0 \frac{d\Phi_{el}}{dt}}_{I_V \text{: Verschiebungsstrom}}
\]

\textbf{Verallgemeinertes Ampèr'sches Gesetz}
\begin{formulabox}
    \[
        \oint_C \Vec{B} \cdot d\Vec{l} = \mu_0 \left( I + I_V \right) = \mu_0 I + \mu_0 \epsilon_0 \frac{d\Phi_{el}}{dt}
    \]
\end{formulabox}

Magnetfelder werden nicht nur von Strömen erzeugt, sondern auch von zetilich variierenden elektrischen Felder. \\[2ex]

\textbf{Gegenstück zum Faraday'schen Gesetz}
\[
    U_{ind} = \oint_C \vec{E} \cdot d\vec{l} = - \frac{d}{dt} \Phi_mag
\]

\subsection{Maxwell-Gleichungen}
\textbf{Gauss'sches Gesetz}
\begin{formulabox}
    \[
        \vec{\nabla} \cdot \vec{E} = \frac{\rho}{\epsilon_0}
    \]
\end{formulabox}

\textbf{Faraday'sches Gesetz}
\begin{formulabox}
    \[
        \vec{\nabla} \times \vec{E} = - \frac{\partial \vec{B}}{\partial t}
    \]
\end{formulabox}

\vspace{5mm}

\textbf{keine magnetische Monopole}
\begin{formulabox}
    \[
        \vec{\nabla} \cdot \vec{B} = 0
    \]
\end{formulabox}

\textbf{Ampèr'sches Gesetz}
\begin{formulabox}
    \[
        \vec{\nabla} \times \vec{B} = \mu_0 \vec{j}  + \mu_0 \epsilon_0 \frac{\partial \vec{E}}{\partial t}
    \]
\end{formulabox}

\subsection{Die Wellengleichung für elektromagnetishe Wellen}
\textbf{Rückblick:} z.B Seilwellen
\[
    \frac{\partial^2 y(x,t)}{\partial x^2} = \frac{1}{v^2}\frac{\partial^2 y(x,t)}{\partial t^2} \hspace{15mm} v^2 = \frac{F_s}{\rho}
\]
Lösungen der Wellengleichung lassen sich Superposition harmonischer Wellenfunktionen schreiben: \\[1ex]
$y(x,t) = y_0 \sin{(kx - \omega t)}$ \hspace{10mm} $k = \frac{2\pi}{\lambda}$: Wellenzahl \\
$y(x,t) = y_0 \sin{(kx + \omega t)}$ \hspace{10mm} $\omega = 2\pi\nu$: Kreisfrequenz \\[2ex]

\textbf{Wellengleichung für elektromagnetische Wellen im Vakuum} \\[1ex]
\hspace{5mm} \textbf{Wellengleichung für elektrisches Feld}
\begin{formulabox}
    \[
        \frac{\partial^2 E}{\partial x^2} = \frac{1}{c^2}\frac{\partial^2 E}{\partial t^2} \hspace{10mm} c^2 = \frac{1}{\mu_0 \epsilon_0}
    \]
\end{formulabox}

\hspace{5mm} \textbf{Wellengleichung für magnetisches Feld}
\begin{formulabox}
    \[
        \frac{\partial^2 B}{\partial x^2} = \frac{1}{c^2}\frac{\partial^2 B}{\partial t^2} \hspace{10mm} c^2 = \frac{1}{\mu_0 \epsilon_0}
    \]
\end{formulabox}

\textbf{Ausbreitungsrichtung} \\
Eine em-Welle breitet sich immer in Richtung von $\vec{E} \times \vec{B}$ aus:

\[
    \vec{E} \times \vec{B} = [\epsilon_0 \sin{(kx - \omega t)} \hat{e}_y] \times [B_0 \sin{(kx - \omega t)} \hat{e}_z]
\]
\begin{formulabox}
    \[
        \Rightarrow E_0 B_0 \sin^2{(kx - \omega t)\hat{e}_x}
    \]
\end{formulabox}
Es gilt: \\
\hspace{3mm} $\rightarrow \vec{E} \perp \vec{B}, \vec{E} \perp \vec{V}$ und $\vec{B} \perp \vec{V}$ \\
\hspace{3mm} $\rightarrow$ em-Wellen sind Transversalwellen \\
\hspace{3mm} $\rightarrow$ Feldstärke: $E = cB$ zu jedem Zeitpunkt an jedem Ort

\subsection{Elektromagnetische Strahlung}
Die \textbf{Energiedichte} einer em-Welle $w_{em}$ setzt sich aus der elektrischen und magnetisches Energiedichte zusammen.\\[1ex]

\hspace{10mm} $\vec{E}-Feld$: $w_{el} = \frac{1}{2}\epsilon_0 E^2$ \\
\hspace{10mm} $\vec{B}-Feld$: $w_{mag} = \frac{1}{2\mu_0} B^2$ \\[1ex]

Für eine em-Welle gilt: $E = cB$
\[
    w_{mag} = \frac{B^2}{2 \mu_0} = \frac{E^2}{2 \mu_0 c^2} = \frac{1}{2}\epsilon_0 E^2 = w_{el}
\]
\textbf{Energiedichte in einer em-Welle}
\begin{formulabox}
    \[
        w_{em} = w_{el} + w_{mag} = \epsilon_0 E^2 = \frac{B^2}{\mu_0} = \frac{EB}{\mu_0 c}
    \]
\end{formulabox}

\textbf{Mittlere Energiedichte} aus Effektivwerten: \\
\[
    E_{eff} = \frac{1}{\sqrt{2}}E_0 \hspace{5mm} \text{und} \hspace{5mm} B_{eff} = \frac{1}{\sqrt{2}}B_0
\]
\textbf{Intensität} (mittlere Leistung pro Fläche senkrecht zur Ausrichtungsrichtung)
\begin{formulabox}
    \[
        I_{em} = \langle w_{em} \rangle \cdot c = \frac{E_{eff} \cdot B_{eff}}{\mu_0} = \frac{1}{2}\frac{E_0 B_0}{\mu_0} = \underbrace{\langle |\vec{s}| \rangle}_{\frac{\vec{E} \times \vec{B}}{\mu_0}}
    \]
\end{formulabox}


    
\end{flushleft}